\begin{tikzpicture}[
  >=Latex,
  line cap=round,
  line join=round,
  font=\small,
  presion/.style={->, orange!85!black, line width=1.15pt},
  radio/.style={->, gray!70!black, line width=0.8pt},
  auxiliar/.style={densely dashed, gray!60, line width=0.8pt}
]
\node[font=\small\bfseries] at (0,2.65) {(a) Corte de la esfera hueca};
\fill[blue!7, even odd rule] (0,0) circle (2.0) (0,0) circle (1.0);
\draw[blue!65!black, line width=1.2pt] (0,0) circle (2.0) (0,0) circle (1.0);
% La presión exterior actúa hacia el centro sobre la superficie r=a.
\foreach \angulo in {0,45,90,135,180,225,270,315} {
  \draw[presion] ({2.55*cos(\angulo)},{2.55*sin(\angulo)})
    -- ({2.05*cos(\angulo)},{2.05*sin(\angulo)});
}
% La presión interna empuja hacia fuera la superficie r=b.
\foreach \angulo in {45,135,225,315} {
  \draw[presion] ({0.52*cos(\angulo)},{0.52*sin(\angulo)})
    -- ({0.98*cos(\angulo)},{0.98*sin(\angulo)});
}
\node[orange!85!black] at (0.15,2.35) {$p_a$};
\node[orange!85!black] at (0,0.3) {$p_b$};
\draw[radio] (0,0) -- (1,0) node[midway,above] {$b$};
\draw[radio] (0,0) -- ({2*cos(-35)},{2*sin(-35)})
  node[pos=0.78,below left] {$a$};
\node[font=\footnotesize,align=center] at (0,-2.45)
  {$0<b<a$: $a$ exterior y $b$ interior};

\begin{scope}[xshift=6.4cm,yshift=0.25cm]
  \node[font=\small\bfseries] at (1.6,2.4) {(b) Condiciones de tracción};
  \draw[->,gray!75] (0,0.55) -- (3.75,0.55) node[right] {$r$};
  \draw[blue!65!black,line width=2pt] (0.55,0.55) -- (3.25,0.55);
  \fill[blue!65!black] (0.55,0.55) circle (2pt) (3.25,0.55) circle (2pt);
  \draw[auxiliar] (0.55,0.55) -- (0.55,1.25);
  \draw[auxiliar] (3.25,0.55) -- (3.25,1.25);
  \node[below] at (0.55,0.45) {$b$};
  \node[below] at (3.25,0.45) {$a$};
  \node[above,align=center] at (0.55,1.3) {$\sigma_{rr}(b)=-p_b$};
  \node[above,align=center] at (3.25,1.3) {$\sigma_{rr}(a)=-p_a$};
  \node[font=\footnotesize,align=center,text width=4cm] at (1.9,-0.35)
    {Intervalo material: $b\le r\le a$\\$0<b<a$};
\end{scope}
\end{tikzpicture}